% !TeX root = main.tex
% Ubah identitas di sini; seluruh halaman awal mengambil data dari berkas ini.
% Nama dosen, NIP/NIDN, dan tanggal masih berupa isian sampai dikonfirmasi.
\newcommand{\judulIndonesia}{Analisis Ketahanan dan Sensitivitas Spektral Representasi Audio dalam Pencarian Rekaman Suara Burung Indonesia dengan Campuran Derau Lapangan ITERA Terkontrol}
\newcommand{\judulInggris}{Analysis of Noise Robustness and Spectral Sensitivity of Audio Representations in Bioacoustic Similarity Retrieval under Controlled ITERA Field-Noise Mixtures in Indonesian Avifauna}
\newcommand{\judulTA}{\MakeUppercase{\judulIndonesia}}
\newcommand{\namaPenulis}{Fabio Banyu Cyto}
\newcommand{\nomorInduk}{123450104}
\newcommand{\namaProgramStudi}{Sains Data}
\newcommand{\namaFakultas}{Sains}
\newcommand{\namaInstitusi}{Institut Teknologi Sumatera}
\newcommand{\namaLokasi}{Lampung Selatan}
\newcommand{\programStudi}{PROGRAM STUDI SAINS DATA}
\newcommand{\fakultas}{FAKULTAS SAINS}
\newcommand{\institusi}{INSTITUT TEKNOLOGI SUMATERA}
\newcommand{\lokasiTA}{LAMPUNG SELATAN}
\newcommand{\tahunTA}{2026}
\newcommand{\pembimbingSatu}{[Nama dan gelar Pembimbing I]}
\newcommand{\identitasPembimbingSatu}{[NIP/NIDN Pembimbing I]}
\newcommand{\pembimbingDua}{[Nama dan gelar Pembimbing II]}
\newcommand{\identitasPembimbingDua}{[NIP/NIDN Pembimbing II]}
\newcommand{\pengujiSatu}{[Nama dan gelar Penguji I]}
\newcommand{\identitasPengujiSatu}{[NIP/NIDN Penguji I]}
\newcommand{\pengujiDua}{[Nama dan gelar Penguji II]}
\newcommand{\identitasPengujiDua}{[NIP/NIDN Penguji II]}
\newcommand{\koordinatorProdi}{[Nama dan gelar Koordinator Program Studi]}
\newcommand{\identitasKoordinator}{[NIP/NIDN Koordinator Program Studi]}
\newcommand{\tanggalPengesahan}{\makebox[3.3cm]{\dotfill}}
\newcommand{\tanggalPernyataan}{\makebox[3.3cm]{\dotfill}}
